% GENERATED FILE. Edit canonical lessons, then run npm run generate:books.
% canonical-source-hash: fnv1a64:d3f6878c15f61651
% canonical-lessons: IT-C272-potrei, IT-C272-sarebbe, IT-C272-ripasso-condizionale, IT-W272-sei-righe

\chapter{Would, Could, Should --- The Condizionale, and Six Lines}
\label{ch:it-the-condizionale}
\hlchaptermodality{272}

\begin{chapteropening}
\textbf{By the end of this chapter:} \emph{I can say what would happen with the condizionale, and write six connected lines about a journey that end looking forward.}

Complete IT-W272-sei-righe: six lines to one reader, scene in the imperfetto, events in order, a reason, and a last line with spero di, sarebbe or se.
\end{chapteropening}

\section[potrei, vorrei, dovrei]{potrei, vorrei, dovrei — the would of Italian}

\label{lesson:IT-C272-potrei}



\begin{quote}
Say \textquotedblleft{}I would like a coffee\textquotedblright{}. (\emph{Vorrei un caffè.}) Say \textquotedblleft{}I should study\textquotedblright{}. (\emph{Dovrei studiare.})

Two forms of one tense, and you have known one of them since chapter eighty-six. This lesson gives the tense its name and its recipe.
\end{quote}

\begin{grammarlens}[title={Grammar Lens: the -ei family}]
This is the \textbf{condizionale}, Italian's \textquotedblleft{}would\textquotedblright{}. Build it on the whole infinitive, minus its final \emph{-e}, and add \textbf{-ei} for \textquotedblleft{}I\textquotedblright{}:

\begin{quote}
parlare, \textbf{parlerei} — I would speak \emph{(the -ar- turns to -er-)} credere, \textbf{crederei} — I would believe dormire, \textbf{dormirei} — I would sleep
\end{quote}

The other persons change only the ending: \textbf{parlerei}, \textbf{parleresti}, \textbf{parlerebbe}, \textbf{parleremmo}, \textbf{parlereste}, \textbf{parlerebbero}.

A few common verbs squeeze their infinitive first, and these are the three you meet most:

\begin{quote}
potere, \textbf{potrei} — I could · volere, \textbf{vorrei} — I'd like · dovere, \textbf{dovrei} — I should
\end{quote}

Three jobs, all of them about softening. A \textbf{polite request}: \emph{Potrei avere un caffè?} — Could I have a coffee? A \textbf{wish}: \emph{Vorrei viaggiare.} And \textbf{advice}: \emph{Al tuo posto, studierei.} — In your place, I'd study.
\end{grammarlens}

\begin{cousinweb}
Latin had no would-tense. Spoken Latin made one by gluing a past of \emph{habēre}, \textquotedblleft{}to have\textquotedblright{}, onto the infinitive: roughly \textquotedblleft{}I had to speak\textquotedblright{}, which slid into \textquotedblleft{}I would speak\textquotedblright{}. Italian chose the perfect, \emph{habuī}, \textquotedblleft{}I had\textquotedblright{}, whose Italian descendant is \emph{ebbi}; worn down on the end of a verb it became \textbf{-ei}, and \emph{ebbe}, \textquotedblleft{}he had\textquotedblright{}, stayed nearly whole in \textbf{parlerebbe}. French and Spanish glued on a different past of the same verb, so their endings do not match. Neither EtymDB nor Lewis \& Short covers a change this late; this is the usual explanation in histories of Romance, given as that. The \emph{habēre} is the same one that became \textbf{avere}.
\end{cousinweb}

\subsection*{Your turn}
Say these aloud:

\begin{itemize}
\raggedright
  \item \emph{parlerei, parleresti, parlerebbe}
  \item the three short ones — \emph{potrei, vorrei, dovrei}
  \item a polite request — \emph{Potrei avere un caffè?}
  \item advice — \emph{Al tuo posto, studierei.}
  \item what you could do now — \emph{Studio l'italiano da tre anni, quindi potrei leggere un romanzo.}
\end{itemize}

\begin{tcolorbox}[breakable,colback=teal!4,colframe=teal!35!black,title={Before you move on}]
Build \textquotedblleft{}I would speak\textquotedblright{}. (\emph{\textbf{Parlerei}} — \emph{-ar-} turns to \emph{-er-}.) What are the three jobs of the condizionale? (A \textbf{polite request}, a \textbf{wish}, \textbf{advice}.) Which Latin verb is hidden in the \emph{-ei}? (\emph{\textbf{Habēre}}, \textquotedblleft{}to have\textquotedblright{} — its past \emph{habuī}, Italian \emph{ebbi}.)
\end{tcolorbox}
\section[sarebbe]{sarebbe — it would be}

\label{lesson:IT-C272-sarebbe}



\begin{quote}
Build \textquotedblleft{}I would speak\textquotedblright{}. (\emph{Parlerei.}) Ask politely for a coffee with \emph{potere}. (\emph{Potrei avere un caffè?})
\end{quote}

\subsection*{You'll want to know: sarebbe}
\textbf{sarebbe} — \textquotedblleft{}it would be\textquotedblright{}, \textquotedblleft{}he or she would be\textquotedblright{}.

\textbf{essere} does not build its condizionale on \emph{esser-}. It uses a stem of its own, \textbf{sar-}, and then the endings are the usual ones: \textbf{sarei}, \textbf{saresti}, \textbf{sarebbe}, \textbf{saremmo}, \textbf{sareste}, \textbf{sarebbero}. \emph{Sarebbe} is the one you will use most, because it judges a situation rather than a person:

\begin{quote}
\emph{Sarebbe difficile.} — It would be difficult. \emph{Sarebbe più comodo in treno.} — It would be more comfortable by train.
\end{quote}

\begin{grammarlens}[title={Grammar Lens: sarebbe bello}]
\textbf{Sarebbe bello} + infinitive — \textquotedblleft{}it would be lovely to\textquotedblright{}:

\begin{quote}
\emph{Sarebbe bello viaggiare insieme.} — It would be lovely to travel together.
\end{quote}

That is a wish with the person left out, softer than \emph{vorrei} and less private than \emph{spero di}. And it sits naturally after a real condition, as the judgement on what follows:

\begin{quote}
\emph{Se piove, sarebbe più comodo prendere l'autobus.} — If it rains, it would be more comfortable to take the bus.
\end{quote}

The \emph{se} half stays in the plain present: it may well rain. The \emph{sarebbe} half is your verdict on the branch.
\end{grammarlens}

\begin{cousinweb}
\textbf{sarebbe} is built on \textbf{essere}, inherited from Latin \emph{esse} (EtymDB 2.1): the stem \emph{sar-} is the one essere uses for what has not happened yet, and \textbf{-ebbe} is the worn Latin \emph{habuit}, \textquotedblleft{}he had\textquotedblright{}, of the last lesson. So \emph{sarebbe} began as something like \textquotedblleft{}he had to be\textquotedblright{}. Neither source records forms this late, and the history of \emph{sar-} itself is not settled, so the stem is described here rather than traced.
\end{cousinweb}

\subsection*{Your turn}
Say these aloud:

\begin{itemize}
\raggedright
  \item \emph{sarei, saresti, sarebbe}
  \item \emph{Sarebbe più comodo in treno.}
  \item \emph{Sarebbe bello viaggiare insieme.}
  \item a verdict with \emph{il più} — \emph{Sarebbe il viaggio più bello del mondo.}
  \item a branch and a verdict — \emph{Se piove, sarebbe più comodo prendere l'autobus.}
\end{itemize}

\begin{tcolorbox}[breakable,colback=teal!4,colframe=teal!35!black,title={Before you move on}]
What stem does \emph{essere} use in the condizionale? (\emph{\textbf{Sar-}}.) Say \textquotedblleft{}it would be lovely to travel together\textquotedblright{}. (\emph{\textbf{Sarebbe bello viaggiare insieme.}}) In \emph{se piove, sarebbe più comodo}, which half is in the plain present? (\emph{\textbf{Se piove.}})
\end{tcolorbox}
\section[(review)]{Review — if this, then that would}

\label{lesson:IT-C272-ripasso-condizionale}



\begin{quote}
Say \textquotedblleft{}it would be difficult\textquotedblright{}. (\emph{Sarebbe difficile.}) Then \textquotedblleft{}I could\textquotedblright{}. (\emph{Potrei.})
\end{quote}

\begin{grammarlens}[title={What you've built}]
Two chapters, and a whole way of talking about what has not happened yet. A real branch: \emph{se piove, resto a casa}. A branch not yet named: \emph{dipende}, or \emph{dipende da te}. And the softened result, in the condizionale: \emph{potrei}, \emph{dovrei}, \emph{sarebbe}.

\begin{quote}
— \emph{Andiamo al mare domani?} — \emph{Dipende. Se c'è il sole, sarebbe bello. Se piove, potremmo andare al museo.}
\end{quote}

\emph{Andiamo} is \textquotedblleft{}shall we go\textquotedblright{}, and \emph{potremmo} is \textquotedblleft{}we could\textquotedblright{}, the \emph{noi} form of \emph{potrei}. Both are the forms you already know how to build, worn by a new person.
\end{grammarlens}

\subsection*{Your turn}
Say these aloud:

\begin{itemize}
\raggedright
  \item \emph{Se c'è il sole, sarebbe bello andare al mare.}
  \item \emph{Se piove, potremmo andare al museo.}
  \item advice inside a branch — \emph{Se sei stanco, dovresti dormire.}
  \item a hope and a verdict — \emph{Spero di viaggiare. Sarebbe bello.}
  \item your own plan for tomorrow, with one \emph{se} and one \emph{sarebbe}
\end{itemize}

\begin{tcolorbox}[breakable,colback=teal!4,colframe=teal!35!black,title={Before you move on}]
In \emph{se piove, potremmo andare al museo}, which half is the condizionale? (\emph{\textbf{Potremmo andare.}}) And which half stays in the plain present? (\emph{\textbf{Se piove.}}) What would \textquotedblleft{}we could\textquotedblright{} be, from \emph{potrei}? (\emph{\textbf{Potremmo.}})
\end{tcolorbox}
\section[(six lines)]{Six lines about a journey}

\label{lesson:IT-W272-sei-righe}



\begin{quote}
Chapter two hundred and sixty-two asked for four lines that held together. This one asks for six, and they tell a story.

Pick one person who will read it: \emph{un amico}, \emph{un'amica}, someone in your family.
\end{quote}

\subsection*{What to know first}
A told story is not a list of events. It has a \textbf{scene}, in the imperfetto: \emph{era}, \emph{c'era}, \emph{faceva}, \emph{ero contento}. It has \textbf{events}, in the passato prossimo, put in order with \emph{prima}, \emph{poi} and \emph{alla fine}. It has at least one \textbf{reason}, pointing back with \emph{perché} or forward with \emph{quindi}.

And a good one ends by looking forward: \emph{spero di …}, or \emph{sarebbe bello …}, or a real condition, \emph{se …}, with the present in the \emph{se} half.

This book's B1 writing paper asks for a connected text of 130 to 170 words. Six lines is the first step towards it, not the whole of it.

\subsection*{Writing — connected composition}
On paper, write \textbf{six lines} to your reader about a journey, real or invented, one line for each of these:

\begin{enumerate}
\raggedright
  \item the scene: when it was, and the weather;
  \item where you were, who was with you, and how long you stayed, with \emph{per};
  \item the first thing that happened, with \emph{prima};
  \item the next thing, with \emph{poi}, and a reason, with \emph{perché} or \emph{quindi};
  \item how it ended and how you felt, with \emph{alla fine} and \emph{ero};
  \item one line forward, with \emph{spero di}, \emph{sarebbe bello} or \emph{se}.
\end{enumerate}

Use only words you already write. The whole exercise is choosing the right past for each line, and joining them.

\begin{tcolorbox}[breakable,colback=teal!4,colframe=teal!35!black,title={Before you move on}]
Now read your lines as you read the journey in chapter two hundred and sixty-four. Mark every verb \textbf{S} for scene or \textbf{E} for event. Lines one and five should be mostly \textbf{S}; three and four mostly \textbf{E}; in line two, a \emph{per} stretch is \textbf{E}. A line about the weather marked \textbf{E} is the commonest slip: change \emph{ha piovuto} back to \emph{pioveva}.

Then two spelling checks. A longer word stressed on its last vowel carries the accent there, \textbf{perché}, \textbf{città}, and so do a few short words you write all the time: \textbf{è}, \textbf{più}, \textbf{sì}. And every double consonant is held: \emph{aspettato}, \emph{abbiamo}, \emph{raccontava}.

Last, read line six alone. If it has a \emph{se}, is the verb after it in the plain present? If it has \emph{sarebbe}, is it your verdict on something that has not happened yet? Then it looks forward, and the story is finished.
\end{tcolorbox}
